\scanpage{5}
\mspage{3}

% section 3, pg 23
\tikz[baseline=(c.base)]{\node[draw,ellipse,inner sep=1pt] (c) {3};}\quad \uline{pg 23}\qquad On peut dire plus simplement que $\,{}_nK$ est étale séparé car $\,{}_nK=\operatorname{Ker}({}_nG\longrightarrow{}_nH)$. En a), ne dirait-tu pas que $\,{}_nK_{\bar s}$ est isomorphe à un sous-groupe de $\,{}_nK_{\bar t}$.

% pg 24
\uline{pg 24}\quad d) $l\,15$ :\quad \tikz[baseline=(c.base)]{\node[draw,ellipse,inner sep=1pt] (c) {$\dim$};}$\operatorname{Ker}(u_{\unsure{s|1}})=0$ ?

% pg 28
\tikz[baseline=(c.base)]{\node[draw,ellipse,inner sep=1pt] (c) {\uline{pg 28}};}\quad Les notions introduites dans 3.5 ont elles été définies pour $S$ non hensélien ? (il faudrait au moins choisir un groupe d'inertie).

% pg 28, below
[item pour 3.7]

% pg 29
\uline{pg 29}\quad $l-2$\quad la terminologie ``clôture normale'' est-elle encore canonique ?

% division
\hrule\smallskip\hrule\smallskip\hrule

% section 4, p 32
\tikz[baseline=(c.base)]{\node[draw,ellipse,inner sep=1pt] (c) {4};}\quad \uline{p 32}\quad $l\,8$\qquad ne dépend \tikz[baseline=(c.base)]{\node[draw,ellipse,inner sep=1pt] (c) {pas};}

% pg 33 ss
\tikz[baseline=(c.base)]{\node[draw,ellipse,inner sep=1pt] (c) {? \quad\uline{pg 33 ss}};}\quad (4.2)\qquad Ne donnerais tu pas un énoncé style ``voisinage tubulaire'' compréhensible. la vérité est

% pg 33 ss, continuation
\deleted{\illegible{1}} on a un morphisme de topos $\operatorname{Spec}(K)_{\mathrm{et}}\xrightarrow{\varphi}\operatorname{Spec}(k)_{\mathrm{et}}$ (``\unsure{ent.|ext.} du corps résiduel'') et $A_0^0$ est un schéma relatif sur $(\operatorname{Spec}(K))_{\mathrm{et}}$ annelé par $\varphi^*\unsure{O|\mathcal O}$ [le $\varphi^*$ des faisceaux d'ensembles] -- Même chose d'ailleurs pour $\underset{\text{ramifier.}}{\varinjlim}\, A_0$

% pg 36
\tikz[baseline=(c.base)]{\node[draw,ellipse,inner sep=1pt] (c) {si ?\quad\uline{pg 36}};}\quad $l-6$.\qquad \unsure{*|``} ce n'est pas la trace qui est contragrédiente ; l'énoncé primitif était correct --

% pg 40
\uline{pg 40}\qquad (4.5.4)\quad Qu'est-il écrit dans la marge ?

% pg 40, l −1
$l-1$\qquad Quel est loc.cit. ?

% pg 40, l 3
$l\,3$\qquad $l$-groupe signifie-t-il tué par $l$, ou par un $l^k$ ?

% last line
le trait est-il strictement local ? ce serait plus agréable.
% archive sheet number: 4
